% !TEX root = ../main.tex
% Wave-10 T3/20: Automorphic correction m(a), tau(a) analysis
% for Delta_5 -> g_{Delta_5} GBKM super-algebra.
% Sources: Borcherds 1995 Invent 109; Gritsenko-Nikulin 1998;
% Lorgat 2020 + slides (/tmp/automorphic-slides.txt).

\providecommand{\LambdaIIone}{\Lambda^{2,1}_{II}}
\providecommand{\PII}{\mathcal P_{II}}
\providecommand{\WII}{W^{(2)}(\LambdaIIone)}
\providecommand{\kappaBKM}{\kappa_{\mathrm{BKM}}}
\providecommand{\kappaCH}{\kappa_{\mathrm{ch}}}
\providecommand{\gDeltaFive}{\mathfrak g_{\Delta_5}}
\providecommand{\ClaimStatusTheorem}{\textsf{[Theorem]}}
\providecommand{\ClaimStatusConjectured}{\textsf{[Conjectured]}}

\section*{Automorphic correction $m(a)$ and $\tau(a)$ for $\gDeltaFive$}

\paragraph{Setup.} On the hyperbolic lattice $\LambdaIIone = \{m f_2 + l f_3
+ n f_{-2} \in \Lambda^{2,1} : m \equiv n \equiv 0 \bmod 2\}$ with Gram
matrix $((\delta_i,\delta_j)) = \bigl(\!\begin{smallmatrix}2&-2&-2\\
-2&-2&2\\ -2&2&-2\end{smallmatrix}\!\bigr)$, Weyl vector
$\rho = \tfrac12(\delta_1+\delta_2+\delta_3) = f_2 - \tfrac12 f_3 + f_{-2}$,
fundamental polyhedron $\PII \subset \mathcal C(\LambdaIIone)_+/\mathbb R_{>0}$
with $\PII{}_{\mathrm{prim}} = \{\delta_1 = 2f_2 - f_3,\ \delta_2 = 2f_{-2} -
f_3,\ \delta_3 = f_3\}$.

\paragraph{The correction coefficient $m(a)$. \ClaimStatusTheorem}
For $a \in \LambdaIIone \cap \mathbb R_{>0}\PII$ and
$(n,l,m) \in \mathbb Z^3$ with $n,m>0$, $n\equiv m \equiv l\equiv 1\bmod 2$,
$4nm - l^2 > 0$ the change-of-variable
$a = (n-1)f_2 - \tfrac{l-1}2 f_3 + (m-1)f_{-2} \in \tfrac12\LambdaIIone$
gives
\begin{equation}\label{eq:m-a-formula}
  m(a) \;=\; -\,\frac{f(n,l,m)}{64},\qquad m(0)=-1,
\end{equation}
where $f(n,l,m) = \sum_{(n,l,m)} f(n,l,m)\,e^{\pi i(nz_1+lz_2+mz_3)}$
are the Fourier coefficients of Igusa's $\Delta_5$. Since $64 \mid f(n,l,m)$
(Maass multiplier quantisation) one has $m(a) \in \mathbb Z$. The sign
$m(a) < 0$ for $(a,a) < 0$ encodes the fermionic imaginary simple roots.

\paragraph{The bosonic exponent $\tau(a_0)$. \ClaimStatusTheorem}
At a null vertex $a_0 \in \{2f_2,\,2f_{-2},\,2f_2 - 2f_3 + 2f_{-2}\}$ of
$\PII$, $(a_0, a_0) = 0$. The Jacobi cusp form identity
$\psi_{5,1/2} = \eta(z_1)^9\nu_{11}(z_1,z_2)$ combined with the
Jacobi-triple-product gives
\begin{equation}\label{eq:tau-9}
  1 \,-\, \sum_{t\in\mathbb N} m(t a_0)\,q^t \;=\;
  \prod_{t\in\mathbb N}(1 - q^t)^{\tau(t a_0)},\qquad \tau(a_0)=9.
\end{equation}
The integer $9$ is the $\eta^9$-exponent and is the bosonic imaginary-null
correction multiplicity.

\paragraph{The $\mathbb Z/2$-graded root system. \ClaimStatusTheorem}
Following Borcherds–Gritsenko–Nikulin, define
\begin{align*}
  \Delta^{\mathrm{re}} &= \PII{}_{\mathrm{prim}}\cdot\mathbb Z
  = \mathbb Z\delta_1 \oplus\mathbb Z\delta_2 \oplus \mathbb Z\delta_3, \\
  \Delta^{\mathrm{im}}_{\bar 0} &= \{\tau(a)\cdot a : (a,a)=0,\ \tau(a)>0\}
  \quad\text{(bosonic null)},\\
  \Delta^{\mathrm{im}}_{\bar 1} &= \{m(a)\cdot a : (a,a)<0,\ m(a)<0\}
  \quad\text{(fermionic negative-norm)},
\end{align*}
so $\Delta = \Delta^{\mathrm{re}} \cup \Delta^{\mathrm{im}}_{\bar 0}
\cup \Delta^{\mathrm{im}}_{\bar 1}$ is a Borcherds–Kac–Moody super-root
system. Generators $\{e_\alpha,h_\alpha,f_\alpha\}_{\alpha\in\Delta}$ with
$\mathfrak{sl}_2$-relations, Serre relations on $\alpha\in\Delta^{\mathrm{re}}$,
and $[\mathfrak g_\alpha,\mathfrak g_{\alpha'}]=0$ whenever $(\alpha,\alpha')=0$
give the GBKM super-algebra $\gDeltaFive$, $\LambdaIIone$-graded with
triangular decomposition $\gDeltaFive = \mathfrak n_-\oplus\mathfrak h
\oplus\mathfrak n_+$, $\dim\mathfrak g_\alpha = \mathrm{mult}_0(\alpha)
- \mathrm{mult}_1(\alpha)$ a supertrace.

\paragraph{Automorphic-corrected denominator. \ClaimStatusTheorem
(Borcherds–Gritsenko–Nikulin)}
\begin{equation}\label{eq:corrected-denominator}
\begin{aligned}
  \tfrac{1}{64}\Delta_5(2z) \;&=\;
  \sum_{w\in\WII}\det(w)\,\Bigl(e^{-2\pi i(w\rho,z)}
  \,-\!\!\!\sum_{a\in\LambdaIIone\cap\mathbb R_{>0}\PII}
  m(a)\,e^{-2\pi i(w(\rho+a),z)}\Bigr) \\
  &=\; e^{-2\pi i(\rho,z)}\prod_{\alpha\in\Delta_+}
  \bigl(1-e^{-2\pi i(\alpha,z)}\bigr)^{\mathrm{mult}\,\alpha}.
\end{aligned}
\end{equation}
The first line is Borcherds' correction: the naive $\sum_w \det(w)
e^{-2\pi i(w\rho,z)}$ fails to be a Siegel modular form; the sum
$-\sum_a m(a) e^{-2\pi i(w(\rho+a),z)}$ over fermionic imaginary simples
with $m(a) = -f(n,l,m)/64$ and the bosonic null exponents $\tau(a_0)=9$
are precisely the corrections that promote the character formula to the
cusp form $\Delta_5(2z)/64$.

\paragraph{Super-multiplicities and $\phi_{0,1}$. \ClaimStatusTheorem}
Setting $\phi_{0,1} = \phi_{12,1}/\delta_{12}
= \sum_{n\geq 0, l\in\mathbb Z} f(n,l)\,e^{2\pi i(nz_1 + lz_2)}$
(weak Jacobi form, weight $0$, index $1$), one has
\begin{equation}\label{eq:mult-phi01}
  \mathrm{mult}\,\alpha \;=\; f(nm,l)\qquad
  \text{for } \alpha = (n,l,m),\; (n,l,m) > 0,
\end{equation}
with $f(0,0)=10$, $f(0,-1)=1$, $f(0,l)=0$ for $l<-1$, and
$f(1,-1)=-64$, $f(1,0)=108$, $f(1,1)=-64$, $f(1,2)=10$ from the
$q^1$-expansion of $\phi_{0,1}$.

\paragraph{First-principles check on $f(1,1,1)=64$. \ClaimStatusTheorem}
From $\phi_{5,1} = \psi_{5,1/2}^2 = \eta(z_1)^{18}\nu_{11}(z_1,z_2)^2$
and the Jacobi-triple expansion $\nu_{11} = q_1^{1/8}q_2^{-1/2}
\prod_{n\geq 1}(1-q_1^{n-1}q_2)(1-q_1^n q_2^{-1})(1-q_1^n)$, the leading
term of $\Delta_5 = \sum \phi_{5,m}(z_1,z_2)e^{\pi i m z_3}$ at
$(n,l,m)=(1,1,1)$ picks out the constant in $\nu_{11}^2 \eta^{18}$ after
the Jacobi-triple expansion, yielding $f(1,1,1) = 64 = 2^6$. The
normalising factor $1/64 = 2^{-6}$ is then forced by $m(0) = -1$ and
integrality of $m(a)$.

\paragraph{Interpretation.} Each $\kappa_{\mathrm{BKM}}(\Phi_N) = c_N(0)/2$
($N \in \{1,2,3,4,6\}$) sits as the \emph{bosonic-null} part of
\eqref{eq:corrected-denominator}; the fermionic part organised by
$m(a) < 0$ is the genuinely super-algebraic correction. At $N=1$,
$c_1(0)/2 = 5 = \kappa_{\mathrm{BKM}}(\Delta_5)$ is the weight of $\Delta_5$, and
the super-algebra $\gDeltaFive$ is the K3$\times E$ BPS quantum group.
The formula $m(a) = -f(n,l,m)/64$ is the universal automorphic-correction
formula propagating through $N \in\{1,2,3,4,6\}$ (Gritsenko–Nikulin 1998).

\paragraph{References.}
Borcherds, \emph{Automorphic forms on $O_{s+2,2}(\mathbb R)$ and infinite
products}, Invent.\ Math.\ 120 (1995) 161--213;
Gritsenko--Nikulin, \emph{Automorphic forms and Lorentzian Kac-Moody
algebras I \& II}, Internat.\ J.\ Math.\ 9 (1998) 153--199, 201--275;
Lorgat, \emph{Automorphic corrections and diagonal divisor Siegel modular
forms of genus 2} (2020).
